% Part: first-order-logic
% Chapter: introduction
% Section: syntax

\documentclass[../../../include/open-logic-section]{subfiles}

\begin{document}

\olfileid{fol}{int}{syn}

\olsection{వాక్యనిర్మాణం}

ఏ సంకేతమాలలను మొదటిస్థాయి విధేయ తర్కపు \tetoken{వాక్యాలుగా}{!!{sentence}s} పరిగణించాలో ముందుగా
ఖచ్చితపరచాలి. దాన్ని తరువాత చేస్తాం; ప్రస్తుతానికి ఉదాహరణల ద్వారా
ముందుకు వెళ్దాం. మొదటిస్థాయి విధేయ తర్కపు ప్రాథమిక నిర్మాణ భాగాలు---దాని
సంకేతావళి---రెండు భాగాలుగా విభజించబడతాయి. నిర్దిష్ట విషయాలను
చెప్పడానికి లేదా నిర్దిష్ట వస్తువులను ఎంచి చూపడానికి ఉపయోగించే
సంకేతాలు మొదటి భాగం. వస్తువులను ఎంచి చూపడానికి \tetoken{స్థిరసంకేతాలను}{!!{constant}s},
ఎంచిన వస్తువుల గురించి ఏదైనా చెప్పడానికి \tetoken{విధేయ సంకేతాలను}{!!{predicate}s} వాడతాం.
ఉదాహరణకు, ఒకే వస్తువును ఎంచి చూపడానికి $\Obj a$ను \tetoken{ఒక స్థిరసంకేతంగా}{!!a{constant}}
వాడి, దాని గురించి
\tetoken{వాక్యం}{!!{sentence}}~$\Atom{\Obj P}{\Obj a}$తో ఏదైనా చెప్పవచ్చు. ``$\Obj
a$'', ``$\Obj P$''లకు అర్థాలను మనసులో పెట్టుకుంటే,
$\Atom{\Obj P}{\Obj a}$ను ఒక సహజభాషా వాక్యంగా చదవవచ్చు (మీరు మొదట
ఔపచారిక తర్కం నేర్చుకున్నప్పుడు బహుశా అలాగే చేసి ఉంటారు). మొదటిస్థాయి
తర్కపు అలాంటి సరళ \tetoken{వాక్యాలు}{!!{sentence}s} లభించిన తరువాత, సంకేతావళి రెండవ
భాగమైన తార్కిక సంకేతాలను (సంయోజకాలు, పరిమాణీకరణికాలు) ఉపయోగించి
మరింత సంక్లిష్టమైన వాటిని నిర్మించవచ్చు. ఉదాహరణకు,
$(\Atom{\Obj P}{\Obj a} \land \Atom{\Obj Q}{\Obj b})$ లేదా
$\lexists[\Obj x][\Atom{\Obj P}{\Obj x}]$ వంటి వ్యక్తీకరణలను
రూపొందించవచ్చు.

కఠినమైన అధితార్కిక అధ్యయనానికి అవసరమైన అర్థవిచార నిర్వచనాలు,
మన \tetoken{వ్యుత్పత్తి}{!!{derivation}} వ్యవస్థల నియమాలను ఇవ్వాలంటే, మొదటిస్థాయి విధేయ తర్కపు
\tetoken{ఒక వాక్యంగా}{!!a{sentence}} దేన్ని పరిగణించాలో ముందుగా ఖచ్చితంగా నిర్వచించాలి.
ప్రాథమిక భావన సులభమే: $\Atom{\Obj P}{\Obj a}$ వంటి కొన్ని సరళ
\tetoken{వాక్యాలను}{!!{sentence}s} కేవలం \tetoken{విధేయ సంకేతాలు}{!!{predicate}s}, \tetoken{స్థిరసంకేతాలు}{!!{constant}s} నుంచి
రూపొందించవచ్చు. వాటి నుంచి సంయోజకాలు, పరిమాణీకరణికాలను ఉపయోగించి
మరింత సంక్లిష్టమైన వాటిని రూపొందిస్తాం. అయితే మరింత సంక్లిష్టమైన
\tetoken{వాక్యాలను}{!!{sentence}s} ఏ నియమాల ప్రకారం రూపొందించడానికి అనుమతి ఉంది?
ఆ నియమాలను తప్పక పేర్కొనాలి; లేకపోతే ``మొదటిస్థాయి విధేయ తర్కపు
\tetoken{వాక్యం}{!!{sentence}}''ను తగినంత ఖచ్చితంగా నిర్వచించినట్లు కాదు. ఇక్కడ కొన్ని
సమస్యలు ఉన్నాయి. సరైన సంకేతమాలలనే \tetoken{వాక్యాలుగా}{!!{sentence}s} లెక్కించడం మొదటి
సమస్య. \emph{అన్ని} \tetoken{వాక్యాలు}{!!{sentence}s} గురించీ గణిత నిరూపణలు ఇవ్వగలిగే
రీతిలో అలా చేయడం రెండవది. చివరగా, ``$!A$లోని ప్రతి~$\Obj x$ స్థానంలో
$\Obj b$ను పెట్టు'' వంటి \tetoken{వాక్యాలపై}{!!{sentence}s} కొన్ని ప్రాథమిక
క్రియలను కూడా ఖచ్చితంగా నిర్వచించాలి. మొదటిస్థాయి విధేయ తర్కంలోని
పరిమాణీకరణికాలు, \tetoken{చరాలు}{!!{variable}s} వల్ల ఇది ఎలా చేయాలో పూర్తిగా
స్పష్టంగా ఉండకపోవడమే ఇబ్బంది. ఉదాహరణకు,
$\lexists[\Obj x][\Atom{\Obj P}{\Obj a}]$ను \tetoken{ఒక వాక్యంగా}{!!a{sentence}}
పరిగణించాలా? $\lexists[\Obj x][\lexists[\Obj x][\Atom{\Obj
P}{\Obj x}]]$ సంగతేమిటి? ``$(\Atom{\Obj P}{\Obj x} \land
\lexists[\Obj x][\Atom{\Obj P}{\Obj x}])$లో $\Obj x$ స్థానంలో
$\Obj b$ను పెట్టు'' అన్నప్పుడు ఫలితం ఏమై ఉండాలి?

\end{document}
